\subsection{Updated Methodology}
\begin{frame}
    \frametitle{Introduce capacity factors to simualtions}
    The current methodology matches installed capacity to a demand 
    in the energy generated. Need to introduce capacity factors to 
    match energy generated to a demand in energy generated.
    \begin{itemize}
        \item Assume \glspl{LWR} operate at a 92.66\% 
              \gls{CF} \cite{us_eia_monthly_2022}
        \item Assume Xe-100 and VOYGR operate at 95\% \gls{CF}
              \cite{xenergy_reactor_nodate,nuscale_technology_nodate}
        \item Assume the \gls{MMR} operates at 100\% \gls{CF}
        \item Do not explicitly model outages for any reactor, extend 
              operating cycle by 1 month for \glspl{LWR} and VOYGRs
        \item This methodology changes the energy supplied by \glspl{LWR}, 
              slightly changes the number of reactors deployed,
              but does not affect the amount of fuel or when the reactors 
              receive fuel. 

    \end{itemize}
\end{frame}

\begin{frame}
    \frametitle{Energy Supplied}
    \begin{columns}
        \column[t]{5cm}
            \begin{itemize}
                \item The demand changes to 87.20 GWe-yr 
                \item All of the scenarios meet the minimum demand
                \item Scenarios 3 (Xe-100) and 6 (Xe-100 + VOYGR) produce 
                      a slight oversupply of energy 
            \end{itemize}
        \column[t]{5cm}
        \vspace{-0.8cm}
            \begin{figure}
                \centering
                \includegraphics[scale=0.4]{updated_energy_after_2025.pdf}
                \caption{Annual average energy produced (GWe-yr) in Scenarios 
                2-7 compared with energy demand}
                \label{fig:updated_energy}

        \end{figure}
    \end{columns}
\end{frame}